% Source: Weinberg GR, physical PDF pp. 61--65
%         (printed pp. 35--39).
% Coverage: Section 2.5, Vectors and Tensors; equations
%           (2.5.1)--(2.5.15).
% Figures/tables/footnotes: none.
% Status: source-reviewed and compile-clean.
% Uncertainties: none.

\subsection{Vectors and Tensors}\label{sec:2.5}

Next we go on to electrodynamics and relativistic hydrodynamics, but it is
convenient first to pause and outline a notation that makes the Lorentz
transformation properties of physical quantities transparent. This notation
will be extended in Chapter~4, on tensor analysis, to encompass general
coordinate transformations, but in fact few changes will be needed.

We have already introduced the term ``four-vector'' for any quantity such as
\(\dd x^\mu\) or \(f^\mu\) or \(p^\mu\) that undergoes the transformation
\begin{equation}
  V^\mu\longrightarrow V'^\mu=\Lambda^\mu{}_\nu V^\nu
  \tag{2.5.1}\label{eq:2.5.1}
\end{equation}
when the coordinate system is transformed by
\begin{equation}
  x^\mu\longrightarrow x'^\mu=\Lambda^\mu{}_\nu x^\nu .
  \tag{2.5.2}\label{eq:2.5.2}
\end{equation}
More precisely, such a \(V^\mu\) should be called a \emph{contravariant}
four-vector, to distinguish it from a \emph{covariant} four-vector, defined as
a quantity \(U_\mu\) whose transformation rule is
\begin{equation}
  U_\mu\longrightarrow U'_\mu=\Lambda_\mu{}^\nu U_\nu ,
  \tag{2.5.3}\label{eq:2.5.3}
\end{equation}
where
\begin{equation}
  \Lambda_\mu{}^\nu
  \equiv\eta_{\mu\rho}\eta^{\nu\sigma}\Lambda^\rho{}_\sigma .
  \tag{2.5.4}\label{eq:2.5.4}
\end{equation}
The matrix \(\eta^{\nu\sigma}\) introduced here is numerically the same as
\(\eta_{\nu\sigma}\), that is,
\begin{equation}
  \eta^{\mu\nu}=\eta_{\mu\nu},
  \tag{2.5.5}\label{eq:2.5.5}
\end{equation}
but we write it with indices upstairs to conform with our summation convention.
Note that
\begin{equation}
  \eta^{\mu\rho}\eta_{\rho\nu}
  =\delta^\mu{}_\nu
  =\begin{cases}
      1,&\mu=\nu,\\
      0,&\mu\ne\nu,
    \end{cases}
  \tag{2.5.6}\label{eq:2.5.6}
\end{equation}
so \(\Lambda_\mu{}^\nu\) is the inverse of the matrix
\(\Lambda^\mu{}_\nu\); that is,
\begin{equation}
  \Lambda_\mu{}^\rho\Lambda^\mu{}_\nu=\delta^\rho{}_\nu,
  \qquad
  \Lambda^\mu{}_\rho\Lambda_\nu{}^\rho=\delta^\mu{}_\nu .
  \tag{2.5.7}\label{eq:2.5.7}
\end{equation}
It follows that the scalar product of a contravariant with a covariant
four-vector is invariant; that is,
\begin{equation}
  U'_\mu V'^\mu
  =\Lambda_\mu{}^\rho\Lambda^\mu{}_\sigma U_\rho V^\sigma
  =U_\rho V^\rho .
  \tag{2.5.8}\label{eq:2.5.8}
\end{equation}
To every contravariant four-vector \(V^\mu\) there corresponds a covariant
four-vector
\begin{equation}
  V_\mu=\eta_{\mu\nu}V^\nu ,
  \tag{2.5.9}\label{eq:2.5.9}
\end{equation}
and to every covariant \(U_\mu\) there corresponds a contravariant
\begin{equation}
  U^\mu=\eta^{\mu\nu}U_\nu .
  \tag{2.5.10}\label{eq:2.5.10}
\end{equation}
Note that raising the index on \(V_\mu\) simply gives back \(V^\mu\), and
lowering the index on \(U^\mu\) simply gives back \(U_\mu\):
\[
  \eta^{\mu\nu}V_\nu
  =\eta^{\mu\nu}\eta_{\nu\rho}V^\rho=V^\mu,
  \qquad
  \eta_{\mu\nu}U^\nu
  =\eta_{\mu\nu}\eta^{\nu\rho}U_\rho=U_\mu .
\]
Note also that Eq.~\eqref{eq:2.5.9} does yield a covariant, because
\[
\begin{aligned}
  V'_\mu
  &=\eta_{\mu\nu}V'^\nu
   =\eta_{\mu\nu}\Lambda^\nu{}_\rho V^\rho\\
  &=\eta_{\mu\nu}\eta^{\rho\sigma}\Lambda^\nu{}_\sigma V_\rho
   =\Lambda_\mu{}^\rho V_\rho ,
\end{aligned}
\]
in agreement with Eq.~\eqref{eq:2.5.3}. Similarly,
Eq.~\eqref{eq:2.5.10} does yield a contravariant.

Although any vector can be written in a contravariant or a covariant form,
there are some vectors, such as \(\dd x^\mu\), that appear more naturally
contravariant and others that appear more naturally covariant. An example of
the latter is the gradient \(\partial_\mu\), which obeys the transformation rule
\[
  \partial_\mu
  \longrightarrow
  \partial'_\mu
  =\frac{\partial x^\nu}{\partial x'^\mu}
   \partial_\nu.
\]
Multiplying Eq.~\eqref{eq:2.5.2} by \(\Lambda_\mu{}^\rho\) gives
\[
  x^\rho=\Lambda_\mu{}^\rho x'^\mu ,
\]
so
\[
  \frac{\partial x^\rho}{\partial x'^\mu}
  =\Lambda_\mu{}^\rho ,
\]
and therefore the gradient is covariant:
\begin{equation}
  \partial'_\mu=\Lambda_\mu{}^\nu\partial_\nu .
  \tag{2.5.11}\label{eq:2.5.11}
\end{equation}
One consequence is that the divergence of a contravariant vector
\(\partial_\mu V^\mu\) is invariant. Another is that the scalar product of
\(\partial_\mu\) with itself, the d'Alembertian operator
\begin{equation}
  \Box\equiv\eta^{\mu\nu}\partial_\mu\partial_\nu
  =\nabla^2-\frac{\partial^2}{\partial t^2},
  \tag{2.5.12}\label{eq:2.5.12}
\end{equation}
is also invariant.

Many physical quantities are not scalars or vectors, but more complicated
objects called \emph{tensors}. A tensor has several contravariant and/or
covariant indices with corresponding Lorentz transformation properties, for
example,
\[
  T^\mu{}_{\nu\rho}\longrightarrow
  T'^\mu{}_{\nu\rho}
  =\Lambda^\mu{}_\sigma\Lambda_\nu{}^\lambda
   \Lambda_\rho{}^\kappa T^\sigma{}_{\lambda\kappa}.
\]
A contravariant or covariant vector can be regarded as a tensor with one index,
and a scalar is a tensor with no indices. There are several ways of forming
tensors out of other tensors:

\noindent(A) \emph{Linear Combinations.} A linear combination of tensors with
the same upper and lower indices is a tensor with these indices. For instance,
if \(R^\mu{}_\nu\) and \(S^\mu{}_\nu\) are tensors, and \(a\) and \(b\) are
scalars, and we define
\[
  T^\mu{}_\nu=aR^\mu{}_\nu+bS^\mu{}_\nu ,
\]
then \(T^\mu{}_\nu\) is a tensor; that is,
\[
\begin{aligned}
  T'^\mu{}_\nu
  &=aR'^\mu{}_\nu+bS'^\mu{}_\nu\\
  &=a\Lambda^\mu{}_\rho\Lambda_\nu{}^\sigma R^\rho{}_\sigma
   +b\Lambda^\mu{}_\rho\Lambda_\nu{}^\sigma S^\rho{}_\sigma\\
  &=\Lambda^\mu{}_\rho\Lambda_\nu{}^\sigma T^\rho{}_\sigma .
\end{aligned}
\]

\noindent(B) \emph{Direct Products.} The product of the components of two
tensors yields a tensor whose upper and lower indices consist of all the upper
and lower indices of the two original tensors. For instance, if
\(A^\mu{}_\nu\) and \(B^\rho\) are tensors, and
\[
  T^{\mu\rho}{}_\nu=A^\mu{}_\nu B^\rho ,
\]
then \(T^{\mu\rho}{}_\nu\) is a tensor; that is,
\[
\begin{aligned}
  T'^{\mu\rho}{}_\nu
  &=A'^\mu{}_\nu B'^\rho\\
  &=\Lambda^\mu{}_\sigma\Lambda^\rho{}_\lambda
    \Lambda_\nu{}^\kappa T^{\sigma\lambda}{}_\kappa .
\end{aligned}
\]

\noindent(C) \emph{Contraction.} Setting an upper and lower index equal and
summing it over its values \(0,1,2,3\) yields a tensor with these two indices
absent. For instance, if \(T^{\mu\rho}{}_{\nu\sigma}\) is a tensor and
\[
  T^\rho{}_\sigma=T^{\mu\rho}{}_{\mu\sigma},
\]
then \(T^\rho{}_\sigma\) is a tensor; that is,
\[
\begin{aligned}
  T'^\rho{}_\sigma
  &=T'^{\mu\rho}{}_{\mu\sigma}\\
  &=\Lambda^\mu{}_\alpha\Lambda^\rho{}_\beta
    \Lambda_\mu{}^\gamma\Lambda_\sigma{}^\delta
    T^{\alpha\beta}{}_{\gamma\delta}\\
  &=\delta^\gamma{}_\alpha\Lambda^\rho{}_\beta
    \Lambda_\sigma{}^\delta T^{\alpha\beta}{}_{\gamma\delta}\\
  &=\Lambda^\rho{}_\beta\Lambda_\sigma{}^\delta
    T^\beta{}_\delta .
\end{aligned}
\]

\noindent(D) \emph{Differentiation.} The derivative \(\partial_\lambda\) of
any tensor is a tensor with one additional lower index \(\lambda\). For
instance, if \(T^{\mu\rho}{}_{\nu\sigma}\) is a tensor and
\[
  T^{\mu\rho}{}_{\nu\sigma\lambda}
  =\partial_\lambda T^{\mu\rho}{}_{\nu\sigma},
\]
then \(T^{\mu\rho}{}_{\nu\sigma\lambda}\) is a tensor; that is,
\[
\begin{aligned}
  T'^{\mu\rho}{}_{\nu\sigma\lambda}
  &=\partial'_\lambda T'^{\mu\rho}{}_{\nu\sigma}\\
  &=\Lambda_\lambda{}^\kappa\partial_\kappa
    \left(
      \Lambda^\mu{}_\alpha\Lambda^\rho{}_\beta
      \Lambda_\nu{}^\gamma\Lambda_\sigma{}^\delta
      T^{\alpha\beta}{}_{\gamma\delta}
    \right)\\
  &=\Lambda^\mu{}_\alpha\Lambda^\rho{}_\beta
    \Lambda_\nu{}^\gamma\Lambda_\sigma{}^\delta
    \Lambda_\lambda{}^\kappa
    T^{\alpha\beta}{}_{\gamma\delta\kappa}.
\end{aligned}
\]
Note that the order of indices matters, even as between upper and lower
indices. For instance, \(T^\mu{}_\nu{}^\rho\) may or may not be the same as
\(T^{\mu\rho}{}_\nu\).

Aside from the scalars, there are three special tensors whose components are
the same in all coordinate systems:

\noindent(i) \emph{The Minkowski Tensor.} The definition of Lorentz
transformations tells us immediately that \(\eta_{\mu\nu}\) is a covariant
tensor:
\[
  \eta_{\mu\nu}
  =\Lambda_\mu{}^\rho\Lambda_\nu{}^\sigma\eta_{\rho\sigma}.
\]
Multiplying this equation by \(\eta^{\mu\alpha}\eta^{\nu\beta}\) and using
Eqs.~\eqref{eq:2.5.6} and \eqref{eq:2.5.4}, we find that
\[
  \eta^{\alpha\beta}
  =\Lambda^\alpha{}_\rho\Lambda^\beta{}_\sigma\eta^{\rho\sigma},
\]
so \(\eta^{\mu\nu}\) is a contravariant tensor. (Recall that
\(\eta_{\mu\nu}\) and \(\eta^{\mu\nu}\) are numerically the same matrix, so
this is a matrix that is both covariant and contravariant.) We can form a mixed
tensor by lowering one index on \(\eta^{\mu\nu}\), or raising one index on
\(\eta_{\mu\nu}\); this gives the Kronecker symbol
\[
  \delta^\mu{}_\nu=\eta^{\mu\rho}\eta_{\rho\nu}.
\]
That this is a tensor follows from rules (B) and (C) and the fact that
\(\eta^{\mu\nu}\) and \(\eta_{\mu\nu}\) are tensors.

\noindent(ii) \emph{The Levi-Civita Tensor.} This is a quantity
\(\varepsilon^{\mu\nu\rho\sigma}\) defined by
\begin{equation}
  \varepsilon^{\mu\nu\rho\sigma}=
  \begin{cases}
    +1,&\mu\nu\rho\sigma\text{ an even permutation of }0123,\\
    -1,&\mu\nu\rho\sigma\text{ an odd permutation of }0123,\\
    0,&\text{otherwise}.
  \end{cases}
  \tag{2.5.13}\label{eq:2.5.13}
\end{equation}
Note that
\[
  \Lambda^\mu{}_\alpha\Lambda^\nu{}_\beta
  \Lambda^\rho{}_\gamma\Lambda^\sigma{}_\delta
  \varepsilon^{\alpha\beta\gamma\delta}
\]
must be proportional to \(\varepsilon^{\mu\nu\rho\sigma}\), because the
left-hand side is odd under any single permutation of the indices
\(\mu,\nu,\rho,\sigma\). To find the constant of proportionality, set
\(\mu\nu\rho\sigma=0123\). The left-hand side is then simply the determinant
of \(\Lambda\), which for proper Lorentz transformations is unity (see
Section~\ref{sec:2.1}). Thus the constant of proportionality is unity; that is,
\begin{equation}
  \Lambda^\mu{}_\alpha\Lambda^\nu{}_\beta
  \Lambda^\rho{}_\gamma\Lambda^\sigma{}_\delta
  \varepsilon^{\alpha\beta\gamma\delta}
  =\varepsilon^{\mu\nu\rho\sigma},
  \tag{2.5.14}\label{eq:2.5.14}
\end{equation}
and therefore \(\varepsilon^{\mu\nu\rho\sigma}\) is a tensor under proper
Lorentz transformations. Under an improper transformation it has the usual
pseudotensor sign.

\noindent(iii) \emph{The Zero Tensor.} We can define a tensor with an
arbitrary pattern of upper and lower indices by setting all its components
equal to zero.

Since \(\eta^{\mu\nu}\) and \(\eta_{\mu\nu}\) are tensors, we can use them to
raise or lower indices on an arbitrary tensor; rules (B) and (C) tell us that
this gives a new tensor with one more upper or lower index and one less lower or
upper index. For instance, if \(T_{\mu\nu\rho}\) is a tensor, then so is
\[
  T_\mu{}^\nu{}_\rho=\eta^{\nu\sigma}T_{\mu\sigma\rho}.
\]
In particular, we can lower some or all of the indices on the Levi-Civita
tensor \(\varepsilon^{\mu\nu\rho\sigma}\). Lowering all the indices gives back
the same numerical quantity except for a minus sign:
\begin{equation}
  \varepsilon_{\mu\nu\rho\sigma}=-\varepsilon^{\mu\nu\rho\sigma}.
  \tag{2.5.15}\label{eq:2.5.15}
\end{equation}

The point of all this algebra is that it enables us to tell at a glance that an
equation is Lorentz-invariant. The fundamental theorem is that \emph{if two
tensors, with the same upper and lower indices, are equal in one coordinate
system, then they are equal in any other coordinate system related to the first
by a Lorentz transformation}. For instance, if
\(T^\mu{}_\nu=S^\mu{}_\nu\), then
\[
  T'^\mu{}_\nu
  =\Lambda^\mu{}_\rho\Lambda_\nu{}^\sigma T^\rho{}_\sigma
  =\Lambda^\mu{}_\rho\Lambda_\nu{}^\sigma S^\rho{}_\sigma
  =S'^\mu{}_\nu .
\]
In particular, the statement that a tensor vanishes is Lorentz-invariant.

The formalism outlined in this section is nothing but a description of the
representations of the homogeneous Lorentz group. We shall explore these
representations in greater generality in Section~\ref{sec:2.12}.
